
\documentclass[9pt,a4paper]{article}
\usepackage[left=1cm,right=1cm,top=1.2cm,bottom=1.2cm]{geometry}
\usepackage[utf8]{inputenc}
\usepackage{multicol}
\usepackage{fancyhdr}
\usepackage{listings}
\usepackage{xcolor}
\usepackage{titlesec}
\usepackage{graphicx}
\graphicspath{{./pics/}}
\usepackage{parskip}
\usepackage{caption}
\usepackage{enumitem}
\usepackage{tikz}
% \usepackage[hidelinks]{hyperref}
\usepackage[]{hyperref}
% \usepackage[utf8]{inputenc}
\usepackage[T1]{fontenc}
\usepackage[ngerman]{babel}
% \usepackage{minted}
\usepackage{float}
\usepackage{amsmath}
\usepackage{wrapfig}
\usepackage{pgfplots}
\usepackage{booktabs}
\usepackage{amssymb}
\usepackage{cancel}
\usepackage{pdfpages}

\usetikzlibrary{positioning}

% Querformat und kompaktes Layout
\pagestyle{fancy}
\fancyhf{}
\fancyfoot[RO]{\thepage}

\setlist[itemize]{nosep,left=0.1cm}  %- kompakte und eingerueckte Listen
\setlist[enumerate]{nosep,left=0.1cm} % aehnlich fuer nummerierte Listen, falls benoetigt

% Abstaende fuer Sektionen nach VHDL-Vorlage [2]
\titlespacing*{\section}{0pt}{0.3em}{0.2em}
\titlespacing*{\subsection}{0pt}{0.2em}{0.1em}
\setlength{\columnsep}{1cm} % - Abstand zwischen Spalten
\setlength{\columnseprule}{0.4pt} % - Linie zwischen Spalten

% Titelformatierung fuer Subsektionen
\titleformat{\subsection}{\bfseries\small}{\thesubsection}{1em}{}

% Farben fuer C++-Syntax-Highlighting
\definecolor{cppkeyword}{rgb}{0.1,0.1,0.7} % Blau-Violett
\definecolor{cpptype}{rgb}{0.4,0.1,0.6} % Violett
\definecolor{cppconcept}{rgb}{0.0,0.5,0.3} % Gruen
\definecolor{cppcomment}{rgb}{0.5,0.5,0.5} % Grau
\definecolor{cppstring}{rgb}{0.7,0.4,0.1} % Orange
\definecolor{cppbackground}{rgb}{0.98, 0.98, 0.98} % Grau
\definecolor{cppframe}{rgb}{0.8,0.8,0.8} % Grau

% Listing-Stil fuer C++
\lstdefinestyle{cppstyle}{
	language=C++, % Sprache C++
	basicstyle=\ttfamily\scriptsize, % - Schriftart und Groesse
	keywordstyle=\color{cppkeyword}\bfseries, % - Schluesselwoerter fett und farbig
	commentstyle=\color{cppcomment}\itshape, % - Kommentare kursiv und farbig
	stringstyle=\color{cppstring}, % - Strings farbig
	backgroundcolor=\color{cppbackground}, % - Hintergrundfarbe
	frame=lines, % - Rahmen
	rulecolor=\color{cppframe}, % - Rahmenfarbe
	breaklines=true, % - Zeilenumbrueche erlauben
	tabsize=2, % - Tabulatorgroesse
	columns=fullflexible, % - Spaltenmodus
	% Weitere Schluesselwoerter und hervorgehobene Typen fuer C++
	morekeywords={ % [4]
			alignas, alignof, and, and_eq, asm, auto, bitand, bitor,
			bool, break, case, catch, char, char16_t, char32_t, class,
			compl, concept, const, constexpr, const_cast, continue,
			decltype, default, delete, do, double, dynamic_cast, else,
			enum, explicit, export, extern, false, float, for, friend,
			goto, if, inline, int, long, mutable, namespace, new, noexcept,
			not, not_eq, nullptr, operator, or, or_eq, private, protected,
			public, register, reinterpret_cast, requires, return, short,
			signed, sizeof, static, static_assert, static_cast, struct,
			switch, template, this, thread_local, throw, true, try,
			typedef, typeid, typename, union, unsigned, using, virtual,
			void, volatile, wchar_t, while, xor, xor_eq, final, override
		},
	emph={ % - Hervorgehobene Typen/Konzepte
			int, double, char, bool, void, std::string, std::vector,
			std::list, std::map, std::set, std::iostream, std::fstream,
			std::strstream, std::iomanip, std::endl, std::cout, std::cin, std::cerr,
			std::exception, std::invalid_argument, std::runtime_error
		},
	emphstyle=\color{cpptype}, % - Stil fuer hervorgehobene Typen
	emph={[2]std::begin, std::end, std::sort, std::for_each,
			malloc, free, new, delete, connect, QObject, QWidget, QPainter,
			typeid, dynamic_cast, static_cast, reinterpret_cast, const_cast}, % - Hervorgehobene Funktionen/Konzeptebene 2
	emphstyle={[2]\color{cppconcept}\bfseries} % - Stil fuer Hervorgehobene Funktionen/Konzeptebene 2
}

\lstset{style=cppstyle}
\newcommand{\sourcecite}[1]{{\small\textsuperscript{[#1]}}} % Fuegt die Zitate als hochgestellte Zahl in Klammern ein
